%-------------------------------------------------
% FileName: chap-1.tex
% Description: 第1章 绪论
%-------------------------------------------------

\clearpage
\setcounter{page}{1}
\pagenumbering{arabic}

\chapter{绪论}

\section{研究背景与意义}

\subsection{研究背景}

近几年来，高校信息化建设速度加快，以前依靠人工记录、打电话排班等手段进行事务管理的事务渐渐迁移到了线上，其中也包含体育馆的日常管理工作。对于高校体育馆这个场景来说，处理起来并不难，只是排场地、收钱、借还器材、统计运营数据而已，不过一遇一起，这些事情就显得零散了——信息只在管理员手里，学生看不到实时状态；同一时间段内有两人的预约出现；月末统计需要翻台账。此类问题需要依靠人工来维持已经变得越来越困难，并且无法给师生提供一个稳定的预约环境。

从使用方的角度来说，普通用户希望随时随地可以打开手机或者电脑查询当前有哪些场地空闲、可以直接选择时段下单、可以看到自己的历史订单；管理员更加关心目前场地是否有人正在使用、器材还有多少、最近的订单和注册情况是怎样的。将这两种需求放在同一个系统中，就是本课题的出发点。

\subsection{研究意义}

本课题的实际意义主要体现在两个比较具体的地方。第一种就是减少体育馆日常管理中重复登记的工作量。以预约为例，如果仍然依靠管理员人工记录，用户每次咨询都要先确认场地、日期和时段，管理员再去核对已有的记录；本系统把场地状态、预约时段、订单状态都存入数据库，用户在页面上即可进行查询与下单，管理员只需处理场地上下架、异常订单、器材归还等少数几项工作。第二是保留业务数据。订单、租借记录、用户、场地信息等都存于表结构当中，以后不管是看高峰时段还是统计器材使用情况，都会比纸质台账容易整理。

从毕业设计训练的角度来说，本系统不是单纯的完成几个页面，而是把一个小型 Web 系统从需求到实现串联起来。开发中会遇到前端路由和页面状态，后端接口分层，数据库表关系，JWT 登录鉴权，订单状态流转和器材库存扣减等问题。所有的内容都和真实项目的基础工程能力相关。系统规模小、没有接入真实的支付、校园统一认证，但是它已经包含了校园资源预约系统中典型的一组业务流程，因此对于类似的场地、教室或者器材管理系统的开发也有一定的参考意义。

\section{国内外研究现状}

\subsection{国外研究现状}

在国外大学、公共体育场馆早就建有线上预约系统。早期系统大多与学校统一身份认证、校园卡或者会员账户绑定在一起，主要解决用户的身份识别、场地的预约、费用的结算以及入馆的核验等问题\cite{feng2026}。平台使用时间越久，系统里面留存的订单以及访问数据也就越来越多，有关的研究也开始重视数据分析方面的技能，比如依照历史预约记录来推测热门时段，依照项目偏好给出场地的推荐，或者借助统计的结果来协助场馆安排开放时间\cite{ding2026}。

国外的相关系统发展大体上是由“线上化预约”到“平台化运营”。它们一般不会只强调某个功能，而更加重视账号体系、支付系统、门禁设备以及数据分析之间的相互配合。但是此类方案的建设成本及维护费用较高，不适合直接套用到本科毕设规模的校园体育馆系统之中。对于本课题而言，更具有参考意义的是它的模块化思想，即预约、订单、用户、权限、统计之间不能混在一起，而应该保持一定的边界。

\subsection{国内研究现状}

国内有关校园公共资源管理系统的相关研究虽然数量不少，但研究对象主要是教室、实验室、图书馆座位、体育场馆等。技术实现上使用的是 Spring Boot、Vue、MySQL 的组合，在论文中会以登录注册、资源查询、预约申请、后台维护、基础统计等为主要功能\cite{zhang2022}。这类系统的开发难度小、数据完整度高，比较适合用作高校毕业设计或者中小型管理系统的技术方案\cite{lin2022}。

但已有研究中也存在着一些容易被忽视的问题。部分系统将重点放在页面是否完整上，而缺少对异常路径的说明，比如重复预约怎样判断、订单取消后状态怎样变化、普通用户能否访问管理接口等等；有些系统的数据统计只是简单地显示数量，并没有和实际的管理决策紧密相连\cite{luo2021}。本课题在实现过程中也受到了以上问题的提醒，在需求分析和设计章节中会单独说明预约冲突校验、角色权限控制、库存扣减和后台数据概览等内容。这样写的不是为了堆叠功能，而是使系统的主体业务流程可以闭环。

\section{研究内容与论文结构}

\subsection{研究内容}

本文主要对体育馆管理系统的开发与实现进行研究，分为三部分。第一层为业务流程梳理，主要是对普通用户和管理员两种角色的操作进行分析，用户侧包含注册登录、浏览场地、选择日期和时段、提交预约、支付或者取消订单、租借器材、查看个人记录；管理侧包含维护场地信息、查看订单、管理用户角色、调整器材库存、处理归还、查看后台概览数据。

第二层就是系统的设计。本文采用前后端分离结构，前端使用 Vue3 和 Element Plus 完成页面与交互，后端使用 Spring Boot、MyBatis-Plus 和 MySQL 完成接口、业务逻辑和数据持久化。数据库的设计以角色、用户、场地、订单、器材、租借记录这六种实体为依托，用状态字段来表现订单支付、取消以及器材归还等业务的变化。

第三层就是功能实现和测试。系统实现了场地分页查询、预约时段展示、订单生成和状态操作、器材租借和归还、后台管理以及数据看板等模块的功能，使用登录认证、预约冲突、库存不足、权限访问等测试用例对主要流程进行了测试。由于本课题时间和场景限制，本系统目前没有接入真实支付、短信提醒、校园统一身份认证等模块，将在以后的展望中予以说明。

\subsection{论文结构安排}

本文一共分为七章。第 1 章为绪论，主要对课题研究背景、研究意义、国内外研究现状、本文的研究内容和论文结构进行介绍；第 2 章为相关的技术、理论基础，介绍系统实现中使用到的各种技术、工具；第 3 章为系统的需求分析，从可行性、功能需求、非功能需求等几个方面说明系统建设的依据。

第四章系统总体设计，主要是对系统架构设计、数据库设计、关键设计方案的论述；第五章系统详细设计与实现，主要是说明各个功能模块的页面实现和业务实现过程；第六章系统测试，对系统运行环境、测试方法、测试结果进行分析；第七章总结与展望，对全文工作进行归纳总结，并提出系统后续完善的方向。
